\section{Related Work}
\paragraph{Pass managers, invalidation, and verifier insertion.}
LLVM's new pass manager caches analyses and uses preservation sets to decide which cached results remain valid after a transformation~\cite{llvmnpm}. Analyses may implement transitive invalidation through their dependencies; passes may also update selected analyses manually. MLIR similarly computes analyses lazily, assumes them invalidated after a pass unless explicitly preserved, and lets an analysis inspect preservation of its dependencies~\cite{mlirpm}. These mechanisms answer whether a cached analysis result may be reused. P1D is the direct demand-driven analogue in our study: an invalidated fact is recomputed when queried.

MLIR also exposes a materially stronger operational alternative: its pass manager can run the IR verifier after each individual pass, and pass instrumentation can insert additional callbacks~\cite{mlirpm}. When one verifier expresses every required relation, this configuration can already provide an effective post-pass gate without obligation metadata. It differs from P2 in three ways: the verifier is ordinarily selected by IR kind rather than by a fact/owner route, it runs after every pass rather than only for due selected relations, and a generic IR verifier cannot enforce a relation absent from that IR's verification interface. Our contribution is therefore not the idea of checking after transformations. It is the fail-closed composition contract that identifies which selected relation became suspect, assigns exactly one production owner, records the earliest artifact boundary at which that route can execute, and carries the obligation there even without a query. We do not report a whole-compilation speed comparison with MLIR-style per-pass verification.

\paragraph{Translation validation.}
Translation validation checks each produced translation~\cite{pnueli98,necula00}; verified validators reduce trust for particular transformations~\cite{tristan08}, and Alive/Alive2 automate bounded LLVM rewrite validation~\cite{alive,alive2}. These define strong relations for specific IRs or transformations. Our complementary problem is scheduling heterogeneous existing verifiers: which selected relation became suspect, who owns it, and when it becomes mandatory. The scheduler neither replaces nor strengthens a routed validator.

\paragraph{Verified, typed, and certifying compilation.}
CompCert proves semantic preservation~\cite{leroy06,leroy09backend}; typed assembly, proof-carrying code, and certifying compilation preserve or check explicit evidence~\cite{tal,pcc,certcompiler}. These provide stronger assurance. Our conditional mechanism merely makes existing production checks unavoidable at selected composition boundaries.

\paragraph{Compiler testing and reduction.}
Csmith, EMI, and YARPGen find compiler failures~\cite{csmith,emi,yarpgen}; C-Reduce, Perses, and delta debugging reduce them~\cite{creduce,perses,delta}. Our targeted faults test scheduling, not natural prevalence or discovery power. Fuzzing can supply future external cases but does not enforce a selected relation before backend consumption.

\begin{table*}[t]
\caption{Feature comparison with the strongest material alternatives. ``Query'' means consumer-triggered execution; ``boundary'' means a declared compiler deadline. This table compares mechanisms, not assurance strength: verified compilation provides a substantially stronger whole-compiler result than our conditional theorem.}
\label{tab:related-comparison}
\scriptsize
\centering
\setlength{\tabcolsep}{2.5pt}
\begin{tabular}{@{}p{0.225\textwidth}cccccccc@{}}
\toprule
Mechanism & Inval. & Sem. facts & Canon. owner & Deadline & Missing route & Selective & S2R gate & Indep. ext. \\
\midrule
LLVM/MLIR analysis managers & yes & analysis & no & query & n/a & query & no & yes \\
Verifier after every pass & no & IR schema & IR owner & each pass & verifier failure & no & yes & yes \\
Generic verifier at chosen sites & no & schema & fixed & call site & n/a & no & no & limited \\
Translation validation / Alive2 & per transform & relation & validator & transform & n/a & per transform & usually no & relation-specific \\
Verified / certifying compiler & proof state & yes & proof/checker & proof boundary & proof failure & proof-guided & yes & proof interface \\
P1D demand recomputation & yes & selected & route owner & query & fail closed on query & query & no mandatory gate & yes \\
P2 obligation-guided & yes & selected & unique & boundary & fail closed & due only & yes & yes \\
P3 always verify & yes & selected & unique & every boundary & fail closed & all routes & yes & yes \\
\bottomrule
\end{tabular}
\end{table*}

Pass managers already invalidate lazily; per-pass verification already gates a complete IR-level checker; validators and verified compilers provide stronger relations or proofs. Our contribution is narrower: map a selected-system invalidation to one owner and deadline, carry it to the earliest executable artifact boundary, and fail closed if no route resolves.

\paragraph{Positioning.}
The strongest practical alternative combines lazy invalidation, production validators, per-pass verification, and testing. Per-pass verification closes the unqueried gap when one available checker expresses the relation. P2 targets heterogeneous or deferred routes whose required artifact or owner appears later. For a monolithic compiler with one complete always-run verifier, P2 adds little.
